\documentclass[../main.tex]{subfiles}

\begin{document}
\subsection{Rental PS Rijal}

\begin{code}
\begin{verbatim}
#include <iostream>
using namespace std;
\end{verbatim}
\end{code}

{\Script} di atas digunakan untuk memberikan kemampuan untuk menggunakan fungsi-fungsi dan objek-objek yang tersedia dalah {\header} {\file} \citecode{iostream}. {\Script} pada baris kedua \citecode{using namespace std;} memberikan kemampuan untuk menggunakan fungsi-fungsi dan objek-objek dari \textit{standard library} tanpa menulis \citecode{std::}.

\begin{code}
\begin{verbatim}
const int user_max = 20;
string users[user_max * 2] = { [0] = "admin", [1] = users[0] };
bool user_membership[user_max];
int user_bills[user_max];
int user_logged_id = -1;
int user_count = 1;
int console_prices[4] = { 0, 5000, 10000, 15000 };
const int food_max = 20;
string foods[food_max];
int food_prices[food_max];
int food_count = 0;
int app_state = 0;
\end{verbatim}
\end{code}

{\Script} di atas digunakan untuk menginisialisasikan semua variabel yang dibutuhkan untuk menjalankan program Rental PS Rijal. Semua variabel yang diawali oleh kata \textit{user} merupakan variabel yang digunakan untuk menyimpan data pengguna atau berkaitan dengan data pengguna. Semua variabel yang diawali oleh kata \textit{food} merupakan variabel yang digunakan untuk menyimpan data makanan dan minuman. Variabel \citecode{app\_state} merupakan variabel yang menyimpan data keadaan aplikasi.

\begin{code}
\begin{verbatim}
void pause() {
  cout << "Press any key to continue. . .\n";
  cin.clear(); cin.ignore(); cin.get(); }
\end{verbatim}
\end{code}

{\Script} di atas merupakan sebuah fungsi tanpa kembalian yang ditandakan oleh kata kunci \citecode{void}. Fungsi tersebut akan menghentikan pengguna sejenak sampai dideteksi sebuah masukan apapun dari \textit{keyboard}.

\begin{code}
\begin{verbatim}
void input_invalid() {
  cin.clear(); cin.ignore();
  cout << "Masukan tidak valid\n"; }
\end{verbatim}
\end{code}

{\Script} di atas merupakan sebuah fungsi tanpa kembalian yang ditandakan oleh kata kunci \citecode{void}. Fungsi tersebut akan memberitahu pengguna bahwa sebuah masukan tidak valid dan membersihkan masukan tidak valid tersebut dari {\ioi} {\stream}.

\begin{code}
\begin{verbatim}
int nselector(int l, int u, string c = ">> ") {
  int n;
  while (true) {
    cout << c;
    cin >> n;
    if (cin.fail() || n < l || n > u) {
      input_invalid();
      continue; } break;
  } return n; }
\end{verbatim}
\end{code}

{\Script} di atas merupakan sebuah fungsi dengan kembalian berupa \citecode{int} atau bilangan bulat. Fungsi tersebut akan menyaring masukan dari pengguna sesuai kedua parameter fungsi \citecode{int l} dan \citecode{int u}. Masukan yang dimasukan harus berada dalam lingkup kedua parameter tersebut untuk dianggap sebagai masukan yang valid.

\begin{code}
\begin{verbatim}
bool login(string name, string password) {
  for (int i = 0; i < user_count; ++i) {
    if (name == users[i * 2]
      && password == users[i * 2 + 1]) {
      user_logged_id = i;
      return true; }
  } return false; }
\end{verbatim}
\end{code}

{\Script} di atas merupakan sebuah fungsi dengan kembalikan berupa tipe data \citecode{bool} yang hanya mampu bernilai bernar atau salah. Fungsi tersebut menerima dua parameter yaitu nama pengguna dan kata sandi. Bila nama pengguna dan kata sandi cocok dengan data yang tersimpan maka keluaran fungsi akan benar dan salah untuk selainnya.

\begin{code}
\begin{verbatim}
bool signin(string name, string password) {
  for (int i = 0; i < user_count; ++i) {
    if (name == users[i * 2]) return false;
  }
  users[user_count * 2] = name;
  users[user_count * 2 + 1] = password;
  user_membership[user_count] = false;
  user_bills[user_count] = 0;
  user_count++;
  return true; }
\end{verbatim}
\end{code}

{\Script} di atas merupakan sebuah fungsi dengan kembalikan berupa tipe data \citecode{bool} yang hanya mampu bernilai bernar atau salah. Fungsi tersebut akan membuat akun pengguna baru bila nama pengguna belum pernah digunakan.

\begin{code}
\begin{verbatim}
void topbar() {
  cout << "================================\n";
  cout << "=========Rental PS Rijal========\n";
  cout << "================================\n\n"; }
\end{verbatim}
\end{code}

{\Script} di atas merupakan sebuah fungsi tanpa kembalian yang ditandakan oleh kata kunci \citecode{void}. Fungsi tersebut digunakan untuk mencetak {\header} atau \textit{bar} pada bagian atas program ketika program dijalankan.

\begin{code}
\begin{verbatim}
void unlogged_page() {
  switch ((app_state & 0x6) >> 0x1) {
    case 0: {
      topbar();
      cout << "1. Login\n2. Sign-In\n";
      cout << "3. Exit\n";
      app_state = (nselector(1, 3) << 0x1);
    } break;
    case 1: {
      cout << "================================\n";
      cout << "=============Log-In=============\n";
      cout << "================================\n\n";
      string name, password;
      cout << "Masukan Username: "; cin >> name;
      cout << "Masukan Password: "; cin >> password;
      if (app_state = login(name, password))
        cout << "Log-In Berhasil\n";
      else cout << "Log-In Tidak Berhasil\n";
      pause();
    } break;
    case 2: {
      cout << "================================\n";
      cout << "=============Sign-In============\n";
      cout << "================================\n\n";
      string name, password;
      cout << "Masukan Username: "; cin >> name;
      cout << "Masukan Password: "; cin >> password;
      if (signin(name, password)) cout << "Sing-In Berhasil\n";
      else cout << "Username sudah ada\n";
      pause(); app_state = 0;
    } break; } }
\end{verbatim}
\end{code}

{\Script} di atas mengandung sebuah fungsi yang digunakan untuk menggambarkan halaman ketika pengguna belum \textit{log-in} atau masuk. Halaman tersebut akan memberikan tiga pilihan yakni \textit{log-in}, \textit{sign-in} dan \textit{exit} yang akan bekerja sesuai namanya.

\begin{code}
\begin{verbatim}
void logged_page() {
  topbar();
  switch ((app_state & 0x6) >> 0x1) {
    case 0: {
      cout << "1. Melakukan Rental PS\n";
      cout << "2. Memesan Makanan\n";
      cout << "3. Daftar Member\n";
      cout << "4. log-out\n";
      app_state |= (nselector(1, 4) << 0x1);
    } break;
\end{verbatim}
\end{code}

{\Script} di atas digunakan untuk menampilkan menu ketika pengguna sudah memasukan akun yang benar dan akun yang dimasukan pengguna bukan merupakan akun admin. Pengguna dapat memilih salah satu pilihan yang tersedia dengan cara memasukan nomor pilihan.

\begin{code}
\begin{verbatim}
    case 1: {
      cout << "Harga Rental PS\n";
      cout << "1. PS 3: 1 Jam (Rp. 5.000)\n";
      cout << "2. PS 4: 1 Jam (Rp. 10.000)\n";
      cout << "3. PS 5: 1 Jam (Rp. 15.000)\n";
      cout << "4. Kembali\n";
      {
        int n = nselector(1, 4);
        if (n == 4) break;
        int m = nselector(1, 9*9*9*9, "Jumlah Jam: ");
        float discount = (float)(console_prices[n] * m) *
          ((0.1f * (m >= 3 && m < 6)) + (0.2f * (m >= 6))
           + (user_membership[user_logged_id] * 0.2f));
        user_bills[user_logged_id]
          += (console_prices[n] * m) - discount;
        pause(); }
      app_state = 1;
    } break;
\end{verbatim}
\end{code}

{\Script} di atas digunakan untuk menangani sistem penyewaan PS. Variabel \citecode{float discount} digunakan untuk menyimpan total diskon yang akan didapatkan oleh pengguna. Variabel \citecode{use\_bills} digunakan untuk menyimpan tagihan pengguna.

\begin{code}
\begin{verbatim}
    case 2: {
      cout << "List Makanan/Minuman\n";
      if (food_count == 0) {
        cout << "Makanan Belum Tersedia\n";
        app_state = 1; pause();
        break; }
      for (int i = 0, j = 0; i < food_count; ++i) {
        cout << i + 1 << ". "
           << foods[i] << " (Rp. "
           << food_prices[i] << ")\n";
      } cout << food_count + 1 << ". Kembali\n";
      int n = nselector(1, food_count + 1);
      if (n == food_count + 1) {
        app_state = 1;
        break; }
      cout << "Makanan yang dipilih: " << foods[n - 1] << "\n";
      cout << "Harga: " << food_prices[n - 1] << "\n";
      user_bills[user_logged_id] += food_prices[n - 1]
            * nselector(1, 9*9*9, "Jumlah: ");
      cout << foods[n - 1] << " berhasil dipesan\n";
      pause();
    } break;
\end{verbatim}
\end{code}

{\Script} di atas digunakan untuk menangani pembelian makanan. Daftar makanan akan ditampilkan sesuai isi variabel \citecode{foods} dan jumlah makanan yang akan ditampilkan sesuai nilai variabel \citecode{food\_count}. Ketika pengguna memesan makanan akan ditambah tagihan berdasarkan harga makanan yang tersedia di \citecode{food\_prices}.

\begin{code}
\begin{verbatim}
    case 3: {
      char choice = 'y';
      while (true) {
        cout << "Apakah anda ingin mendaftar sebagai"
           << " member rental PS?(Y/T): ";
        cin >> choice;
        if (cin.fail() ||
          ((choice != 'Y' && choice != 'y') &&
           (choice != 'T' && choice != 't'))) {
          input_invalid();
          continue;
        } break; }
      if (choice == 'Y' || choice == 'y') {
        if (user_membership[user_logged_id])
          cout << "Anda sudah terdaftar sebagai member\n";
        else {
          cout << "Anda berhasil mendaftar sebagai member\n";
          user_membership[user_logged_id] = true; }
      } app_state = 1; pause();
    } break; } }
\end{verbatim}
\end{code}

{\Script} di atas menangani pendaftaran \textit{membership} akun pengguna. Bila akun pengguna yang memiliki \textit{membership} maka nilai yang disimpan pada variabel {\sarray} \citecode{user\_membersip} untuk akun pengguna saat ini akan bernilai \citecode{true}.

\begin{code}
\begin{verbatim}
void admin_page() {
  topbar();
  switch ((app_state & 0x6) >> 1) {
    case 0: {
      cout << "1. Jual Makanan\n";
      cout << "2. Data Penyewa\n";
      cout << "3. log-out\n";
      app_state |= (nselector(1, 3) << 0x1);
    } break;
\end{verbatim}
\end{code}

{\Script} di atas digunakan untuk menampilkan pilihan-pilhan yang tersedia pada menu admin. Pada menu admin tersedia tiga pilihan yakni jual makanan, data penyewa dan \textit{log-out}. Pengguna dapat memilih salah satu dari tiga pilihan tersebut dengan cara memasukan nomor pilihan yang ditampilkan.

\begin{code}
\begin{verbatim}
    case 1: {
      cout << "List Makanan/Minuman\n";
      if (food_count == 0) {
        cout << "Makanan/Minuman Belum Tersedia\n";
      } else {
        for (int i = 0; i < food_count; ++i) {
          if (app_state >> 4 == 0) cout << "~ ";
          else cout << i + 1 << ". ";
          cout << foods[i] << " (Rp. "
             << food_prices[i] << ")\n";
        }
        if (app_state >> 4 != 0)
          cout << food_count + 1 << ". Kembali\n";
        else cout << "\n";
      }
\end{verbatim}
\end{code}

{\Script} di atas digunakan untuk menampilkan daftar makanan yang tersimpan pada variabel {\sarray} \citecode{foods}. Jumlah makanan yang akan ditampilkan sesuai dengan nilai variabel \citecode{food\_count} yang bertipe data \citecode{int}.

\begin{code}
\begin{verbatim}
      if ((app_state >> 4) == 0) {
        cout << "1. Masukan Makanan\n";
        cout << "2. Hapus Makanan\n";
        cout << "3. Kembali\n";
        int n = nselector(1, 3);
        if (n == 1) {
          cout << "\33[H\33[2J"; topbar();
          cout << "Masukan Makanan/Minuman: ";
          cin.clear(); cin.ignore();
          getline(cin, foods[food_count]);
          cout << "Masukan Harga: ";
          cin >> food_prices[food_count];
          cout << "Makanan/Minuman Berhasil"
             << " Dimasukan\n";
          food_count++;
          pause(); app_state = 3; break;
        } else if (n == 3) {
          app_state = 1; break;
        }
        app_state |= (0x1 << 4);
        break;
      }
      int n = nselector(1, food_count + 1);
      if (n == food_count + 1) {
        app_state = 3;
        break;
      }
      cout << foods[n - 1] << " berhasil dihapus\n";
      for (int i = (n - 1); i < food_max - 1; ++i) {
        foods[i] = foods[i + 1];
        food_prices[i] = food_prices[i + 1];
      } food_count--;
      pause(); app_state = 3;
    } break;
\end{verbatim}
\end{code}

{\Script} di atas digunakan untuk menghapus dan menambah {\entry} makanan atau minuman dalam variabel {\sarray} \citecode{foods}. Bila makanan ditambah maka nilai \citecode{food\_count} akan dinaikan dan data makanan tersebut akan dimasukan ke indeks sesuai nilai \citecode{food\_count} pada \citecode{foods}. Penghapusan makanan dilakukan sebaliknya yakni mengurangi nilai \citecode{food\_count} dan menggeser nilai \citecode{foods} hingga tidak tersedia nilai yang kosong.

\begin{code}
\begin{verbatim}
    case 2: {
      cout << "Data Penyewa\n";
      for (int i = 1; i < user_count; ++i) {
        cout << i << ". " << "Username: "
           << users[i * 2] << " (Rp."
           << user_bills[i] << ")\n";
      }
      app_state = 1; pause();
    } break; } }
\end{verbatim}
\end{code}

{\Script} di atas digunakan untuk menampilkan data penyewa berupa daftar nama-nama pengguna yang tersedia pada variabel {\sarray} \citecode{users} dan tagihan pengguna yang tersedia pada variabel {\sarray} \citecode{user\_bills}.

\begin{code}
\begin{verbatim}
int main() {
  while (true) {
    cout << "\33[H\33[2J";
    if (app_state == 0x6) return 0;
    else if (app_state == 0x9) {
      user_logged_id = -1;
      app_state = 0; }
    else if (user_logged_id == 0 &&
         app_state == 0x7) {
      user_logged_id = -1;
      app_state = 0; }
    switch (app_state & 0x1) {
      case 0: {
        unlogged_page();
      } break;
      case 1: {
        if (user_logged_id == 0) admin_page();
        else logged_page();
      } } } }
\end{verbatim}
\end{code}

{\Script} di atas merupakan \textit{entry point} dari program Rental PS Rijal. Fungsi tersebut akan menjalankan fungsi-fungsi lainnya tergantung dengan nilai \citecode{app\_state}. Program tersebut akan dijalankan selamanya, selama nilai \citecode{app\_state} bukan 9.

\subsection{Permasalahan Matematika Diskrit dan Aljabar Linier}

\begin{code}
\begin{verbatim}
int factorial(int n, int acc = 1)
{
  if (n == 1) return acc;
  return factorial(n - 1, acc * n);
}
\end{verbatim}
\end{code}

{\Script} di atas mengandung sebuah fungsi yang mengembalikan nilai bilangan bulat dengan mengambil masukan berupa \citecode{int n}. Fungsi tersebut akan mengembalikan sebuah nilai faktorial dari parameter yang dimasukan dengan menggunakan metode rekursif.

\begin{code}
\begin{verbatim}
cout << "1. Permutasi\n";
cout << "2. Kombinasi\n";
cout << "3. SPL\n";
int choice;
while (true) {
  cout << "Pilih: ";
  cin >> choice;
  if (cin.fail() || choice < 1
    || choice > 3) {
    cin.clear();
    cin.ignore();
    cout << "Masukan tidak valid." << endl;
    continue;
  } break;
}
\end{verbatim}
\end{code}

{\Script} di atas akan mencetak tiga pilihan yang diberikan kepada pengguna. Ketiga pilihan tersebut ialah permutasi, kombinasi dan sistem persamaan linier. Masukan akan diminta kemudian akan dicek apakah masukan lebih kecil dari 1 atau lebih besar dari 3.

\begin{code}
\begin{verbatim}
case 1: {
  int n, r;
  cout << "Masukan n: ";
  cin >> n;
  cout << "Masukan r: ";
  cin >> r;
  cout << "Hasil Permutasi: " 
     << factorial(n)/factorial(n - r)
     << endl;
} break;
\end{verbatim}
\end{code}

{\Script} di atas digunakan untuk menghasilkan nilai permutasi dari dua masukan. Nilai permutasi dicari dengan penggunakan fungsi \citecode{factorial} dan meletakannya sesuai rumus permutasi yang mengambil masukan \citecode{n} dan \citecode{r}.

\begin{code}
\begin{verbatim}
case 2: {
  int n, r;
  cout << "Masukan n: ";
  cin >> n;
  cout << "Masukan r: ";
  cin >> r;
  cout << "Hasil Kombinasi: " 
     << factorial(n)/(factorial(n - r) * factorial(r))
     << endl;
} break;
\end{verbatim}
\end{code}

{\Script} di atas digunakan untuk menghasilkan nilai kombinasi dari dua masukan. Nilai kombinasi dicari dengan penggunakan fungsi \citecode{factorial} dan meletakannya sesuai rumus kombinasi sesuai masukan \citecode{n} dan \citecode{r}.

\begin{code}
\begin{verbatim}
case 3: {
  cout << "Masukan Nilai Matriks 3x3\n";
  float matrix[12];
  for (int i = 0; i < 3; ++i) {
    for (int j = 0, mid = 0; j < 3; ++j) {
      cin >> matrix[i * 4 + j]; } }
\end{verbatim}
\end{code}

{\Script} di atas digunakan untuk meminta masukan matriks 3x3 kepada pengguna. Masukan tersebut kemudian akan disimpan ke dalam variabel {\sarray} \citecode{float matrix[12]}. Perulangan akan dilakukan selama 9 kali untuk mengisi nilai variabel {\sarray} \citecode{matrix}.

\begin{code}
\begin{verbatim}
  cout << "Masukan Nilai B\n";
  for (int i = 0; i < 3; ++i) {
    cin >> matrix[(i + 1) * 4 - 1];
  } cout << "Augmented Matriks\n";
  for (int i = 0; i < 3; ++i) {
    for (int j = 0; j < 4; ++j) {
      cout << matrix[i * 4 + j] << " ";
    } cout << endl;
  }
\end{verbatim}
\end{code}

{\Script} di atas digunakan untuk meminta nilai B yang berupa vektor ruang 3 dimensi. Nilai vektor B tersebut akan digabungkan ke variabel {\sarray} \citecode{matrix} kemudian nilai yang sudah disimpan tersebut akan ditampilkan kembali kepada pengguna.

\begin{code}
\begin{verbatim}
  float x;
  if (matrix[4] != 0 && matrix[8] != 0
    && matrix[9] != 0) {
    x = matrix[4]/matrix[0];
    for (int i = 0; i < 4; ++i)
      matrix[4 + i] += matrix[i] * -x;
    x = matrix[8]/matrix[0];
    for (int i = 0; i < 4; ++i)
      matrix[8 + i] += matrix[i] * -x;
    x = matrix[9]/matrix[5];
    for (int i = 0; i < 4; ++i)
      matrix[8 + i] += matrix[4 + i] * -x;
  }
  cout << "Matriks Segitiga Atas | SPL\n";
  for (int i = 0; i < 3; ++i) {
    for (int j = 0; j < 3; ++j) {
      cout << matrix[i * 4 + j] << " ";
    } cout << " | " << matrix[(i + 1) * 4 - 1] << endl;
  }
\end{verbatim}
\end{code}

{\Script} di atas digunakan untuk membuat matriks segitiga atas dengan cara membagi bagian segitiga bawah dari \citecode{matrix} menjadi 0. Nilai yang digunakan untuk membaginya disimpan pada variabel \citecode{x}.

\begin{code}
\begin{verbatim}
  float z = matrix[11]/matrix[10];
  float y = (matrix[7] - (matrix[6] * z))/matrix[5];
  x = (matrix[3] - (matrix[2] * z) - matrix[1] * y)/matrix[0];
  cout << "x: " << x << "\n";
  cout << "y: " << y << "\n";
  cout << "z: " << z << "\n";
} break; } }
\end{verbatim}
\end{code}

{\Script} di atas digunakan untuk menghasilkan sistem persamaan linier yang dihasilkan dari masukan berupa matriks 3 kali 3 dan vektor B yang digabungkan menjadi satu buah matriks 3 kali 4. Nilai sistem persamaan linier yang telah diproses akan ditampilkan pada akhirnya kepada pengguna.
\end{document}
